\begin{abstract}
Public chest X-ray (CXR) collections are often merged into multi-class benchmarks in which
patients, labels and image sources are not controlled, so internal accuracy can overstate how a
model behaves at a new hospital. We rebuild a six-class benchmark (COVID-19, Normal, Pleural
Effusion, Pneumonia, Pulmonary Nodule/Mass, Tuberculosis) from official sources, including TBX11K and the
COVID-19 Radiography Database, with explicit label rules, per-source caps, patient-grouped splits (patient identity is unknown for two
sources) and duplicate checks. An earlier version of this pipeline reported \RVFiveAcc\% internal accuracy, but an audit
found patient leakage, proxy labels and an external specificity of \RVFiveExtSpecK/\RVFiveExtSpecN. We
train an ImageNet-initialised DenseNet-121 (seed \RSeed, CPU) and evaluate the test split and an external
Montgomery County cohort once. On the leakage-controlled test split the model reaches \RTestAcc\% accuracy
(macro-F1 \RTestMacroFOne\%, macro ROC-AUC \RTestMacroAUC), but on the external cohort it labels only
\RExtSpecK/\RExtSpecN{} normal scans as Normal and names \RExtTBK/\RExtTBN{} tuberculosis scans as
Tuberculosis. A first training run called \RRunOneExtCovidCalls{} of \RNExternal{} external scans COVID-19; with a second build,
which added TBX11K and the COVID-19 Radiography Database among other changes, only \RExtCovidCalls{} were, but the errors moved
to other abnormal classes and the external gap remained. This is consistent with, but does not prove, a source
shortcut (indirect evidence). Results come from a single run in a
retrospective, public-data setting and are not evidence of clinical validity.
\end{abstract}
